\section{RQ3: Controlling B-cos Explanations}
\label{sec:rq3_results}

RQ3 asks whether spatial supervision can control the attribution distribution of a
learned policy without substantially reducing task performance. Five independently
trained $B=2.3$ PPO policies from RQ1 were used as source checkpoints. Each checkpoint
was continued for 303,104 transitions under four matched conditions:

\begin{itemize}
    \item pure PPO continuation without explanation supervision;
    \item supervision using the strict sequential task-dependent mask;
    \item supervision using a local action-conditioned mask; and
    \item supervision using the exactly flipped strict sequential mask.
\end{itemize}

For movement actions, the local mask selects the agent and the single
action-relative cell; for interaction actions, it selects the agent and a valid
front-cell affordance. All remaining pixels are penalized.

The supervised branches used an explanation coefficient of $0.15$ and limited the
explanation actor-gradient norm to at most $0.1$ times the PPO actor-gradient norm.
For an RGB6 explanation $E\in\mathbb{R}^{6\times H\times W}$, attribution energy was
computed as
\begin{equation}
    U(x,y)=\sum_{c=1}^{6}\lvert E_c(x,y)\rvert.
    \label{eq:rq3_channelwise_energy}
\end{equation}
Thus, opposing contributions in different channels cannot cancel. The localization
score is the fraction of total $U$ inside the positive region of the evaluated mask.

Every final policy was evaluated over 100 stochastic episodes. Explanation metrics
used the same 256 fixed state--action samples for all branches within each seed. Each
branch was evaluated under the sequential, local, and flipped masks, rather than only
under its training mask. Reported intervals are 95\% bootstrap confidence intervals
over the five independent seeds.

\begin{table}[htbp]
    \centering
    \small
    \setlength{\tabcolsep}{4pt}
    \resizebox{\textwidth}{!}{%
    \begin{tabular}{lccccc}
        \toprule
        Branch & Success & Return & Episode length & Deletion advantage & Insertion advantage \\
        \midrule
        Pure PPO
        & $1.000 \pm 0.000$
        & $1.125 \pm 0.0004$
        & $17.71 \pm 0.30$
        & $0.252 \pm 0.078$
        & $0.292 \pm 0.122$ \\
        Sequential
        & $0.998 \pm 0.003$
        & $1.123 \pm 0.004$
        & $19.37 \pm 2.73$
        & $0.451 \pm 0.039$
        & $0.241 \pm 0.098$ \\
        Local action
        & $0.988 \pm 0.018$
        & $1.112 \pm 0.018$
        & $25.96 \pm 11.62$
        & $0.230 \pm 0.194$
        & $0.172 \pm 0.069$ \\
        Flipped
        & $0.994 \pm 0.009$
        & $1.115 \pm 0.016$
        & $24.66 \pm 11.01$
        & $0.219 \pm 0.056$
        & $0.211 \pm 0.070$ \\
        \bottomrule
    \end{tabular}%
    }
    \caption{Policy quality and perturbation faithfulness over five matched seeds.
    Values are means $\pm$ half-widths of 95\% bootstrap confidence intervals. Deletion and insertion
    are reported as advantages over equal-area random perturbations; higher is better.}
    \label{tab:rq3_results}
\end{table}

Task performance was largely preserved (Table~\ref{tab:rq3_results}). Mean success was
$1.000$ for pure continuation, $0.998$ for sequential supervision, $0.988$ for the
local mask, and $0.994$ for the flipped control. Their paired success differences
relative to pure were $-0.002$ $[-0.006,0.000]$, $-0.012$
$[-0.036,0.000]$, and $-0.006$ $[-0.018,0.000]$, respectively.

\begin{table}[htbp]
    \centering
    \small
    \resizebox{\textwidth}{!}{%
    \begin{tabular}{lccc}
        \toprule
        Branch & Sequential-mask ratio & Local-mask ratio & Flipped-mask ratio \\
        \midrule
        Pure PPO   & $0.252 \pm 0.038$ & $0.190 \pm 0.035$ & $0.748 \pm 0.038$ \\
        Sequential & $0.644 \pm 0.083$ & $0.394 \pm 0.078$ & $0.356 \pm 0.081$ \\
        Local action & $0.546 \pm 0.076$ & $0.480 \pm 0.100$ & $0.454 \pm 0.076$ \\
        Flipped    & $0.090 \pm 0.024$ & $0.076 \pm 0.018$ & $0.910 \pm 0.023$ \\
        \bottomrule
    \end{tabular}%
    }
    \caption{Channelwise absolute attribution-energy ratios under all three spatial masks.
    Values are means $\pm$ half-widths of 95\% bootstrap confidence intervals over
    five matched seeds.}
    \label{tab:rq3_localization}
\end{table}

Sequential supervision increased sequential-mask alignment from $0.252$ to $0.644$,
a paired increase of $0.392$ $[0.305,0.468]$. Local supervision increased local-mask
alignment from $0.190$ to $0.480$, a paired increase of $0.290$
$[0.208,0.365]$. The flipped control increased flipped-mask alignment from $0.748$
to $0.910$, a paired increase of $0.162$ $[0.130,0.188]$, while sequential-mask
alignment fell to $0.090$. Thus, optimization follows the supplied spatial target
rather than merely concentrating an already task-aligned explanation.

The initial loss instead formed
$U_{\mathrm{old}}(x,y)=\lvert\sum_c E_c(x,y)\rvert$. This was unsuitable for
RGB6 because oppositely signed contributions from complementary channels can cancel
before spatial localization is measured. In the example in
Figure~\ref{fig:rq3_old_vs_channelwise}, supporting and suppressing contributions
occur across the six channels at the same spatial locations. The final definition,
$U(x,y)=\sum_c\lvert E_c(x,y)\rvert$, takes the magnitude before summing channels,
so these contributions cannot cancel inside the mask loss.

\begin{figure}[htbp]
    \centering
    \includegraphics[width=\textwidth]{Thesis_Template_DE_EN/BA_MA_English/sections/images/rq3_old_vs_channelwise_loss_rgb6_seed2_state27.pdf}
    \caption{Signed RGB6 explanations for pure PPO, the initial loss, the final
    channelwise loss, and the local-action-mask loss on state~27 and the same toggle action. Red supports the selected
    action, blue suppresses it, and white is zero. The last column is the signed sum of all six
    channels. Each row is normalized independently, while one symmetric scale is
    shared by the six channels and their sum; the color bars report the original
    contributions. Opposite signs across channels show why
    summing before taking absolute values permits cancellation. The
    trained checkpoints differ in continuation budget and mask variant; the comparison
    is therefore illustrative rather than quantitative.}
    \label{fig:rq3_old_vs_channelwise}
\end{figure}

\begin{figure}[htbp]
    \centering
    \begin{subfigure}[t]{0.98\textwidth}
        \centering
        \includegraphics[width=\textwidth]{Thesis_Template_DE_EN/BA_MA_English/sections/images/rq3_rgb6_channels_state_10_toggle_seed2_compacted_absolute_energy.pdf}
        \caption{Toggle action.}
        \label{fig:rq3_toggle_explanation}
    \end{subfigure}

    \vspace{0.8em}

    \begin{subfigure}[t]{0.98\textwidth}
        \centering
        \includegraphics[width=\textwidth]{Thesis_Template_DE_EN/BA_MA_English/sections/images/rq3_rgb6_channels_state_22_pickup_seed2_compacted_absolute_energy.pdf}
        \caption{Pickup action.}
        \label{fig:rq3_pickup_explanation}
    \end{subfigure}

    \caption{Representative RGB6 explanations from seed~2. Each row compares the
    observation, pure PPO, sequentially supervised PPO, and PPO supervised with the
    local action mask. The six channels are
    compacted using $U=\sum_c\lvert E_c\rvert$. Each explanation is normalized
    independently on a linear scale. White denotes zero, magenta denotes higher
    attribution energy, and the color bars report the original magnitudes.}
    \label{fig:rq3_representative_explanations}
\end{figure}

\begin{figure}[htbp]
    \centering
    \begin{subfigure}[t]{0.98\textwidth}
        \centering
        \includegraphics[width=\textwidth]{Thesis_Template_DE_EN/BA_MA_English/sections/images/rq3_rgb6_channels_state_10_toggle_seed2_absolute.pdf}
        \caption{Toggle action.}
        \label{fig:rq3_toggle_six_channels}
    \end{subfigure}

    \vspace{0.8em}

    \begin{subfigure}[t]{0.98\textwidth}
        \centering
        \includegraphics[width=\textwidth]{Thesis_Template_DE_EN/BA_MA_English/sections/images/rq3_rgb6_channels_state_22_pickup_seed2_absolute.pdf}
        \caption{Pickup action.}
        \label{fig:rq3_pickup_six_channels}
    \end{subfigure}

    \caption{The same representative explanations before channel compaction. Rows
    correspond to pure PPO, sequentially supervised PPO, and PPO supervised with the
    local action mask; columns show the RGB6
    channels $R$, $G$, $B$, $255-R$, $255-G$, and $255-B$. Absolute attribution
    magnitude is shown from white (zero) to magenta. For each model and observation,
    one linear scale is shared across all six channels, and the color bar reports the
    original magnitude.}
    \label{fig:rq3_six_channel_explanations}
\end{figure}

Localization and perturbation faithfulness were related but not equivalent.
Sequential supervision improved deletion advantage over pure PPO by $0.199$
$[0.098,0.290]$, whereas its insertion difference was $-0.051$
$[-0.191,0.060]$. For the local branch, deletion did not change conclusively and
insertion decreased by $0.120$ $[-0.247,-0.031]$. Neither faithfulness difference
was conclusive for the flipped branch. Hence, spatial compliance did not imply a
general improvement across all faithfulness criteria.

\subsection{Answer to RQ3}

Channelwise absolute-energy supervision reliably changed where pretrained B-cos PPO
policies attributed their selected actions while largely preserving DoorKey task
performance. Sequential and local action-conditioned supervision moved attribution
energy toward their intended regions, and the flipped negative control moved it in
the opposite direction. Explanation control therefore follows the supplied mask, but high mask
alignment alone does not guarantee semantic correctness or uniformly improved
faithfulness.
